Para asegurar el correcto funcionamiento de la aplicación, se han desarrollado pruebas tanto en el backend como en el frontend.

\paragraph{TestCase}
Las pruebas del backend se han implementado utilizando la librería TestCase \cite{djangotesting2025}. Se han creado pruebas específicos para cada uno de los módulos principales del sistema:

\begin{description}
    \item Módulo de autenticación: se verifican el registro, inicio de sesión y verificación del token.
    \item Módulo de amistad: se comprueba el envío, aceptación y eliminación de solicitudes de amistad.
    \item Módulo de búsqueda: se validan las consultas de videojuegos y listas.
\end{description}

Estas pruebas se ejecutan utilizando una base de datos temporal que se crea y destruye para cada prueba, lo que permite no afectar a los datos reales del sistema. 
Cada prueba se ejecuta siguiendo un patrón:
\begin{enumerate}
    \item Preparación del entorno: se crea un usuario de prueba y se insertan videojuegos de prueba.
    \item Ejecución de la petición: se realiza una petición HTTP al endpoint correspondiente del servidor.
    \item Verificación de la respuesta: se comprueba el código de estado HTTP y el contenido de la respuesta utilizando aserciones.
\end{enumerate} 

En el código \ref{COD:CODIGOPYE} se muestra un ejemplo de prueba de búsqueda de videojuegos y en la figura \ref{FIG:TESTCASE} se muestra su salida.
\PythonCode[COD:CODIGOPYE]{Pruebas django}{En este código se presenta la ejecución y verificación de la respuesta}{python/tests.py}{95}{99}{1}

\begin{figure}[Salida TestCase]{FIG:TESTCASE}{Salida de una prueba de TestCase que comprueba la búsqueda de videojuegos}
    \image{10cm}{}{django}
\end{figure}

\paragraph{Cypress}
Para el frontend, se han realizado pruebas end-to-end utilizando la herramienta Cypress \cite{cypresstest2025}, que permite simular la interacción del usuario con la interfaz.

Se han desarrollado tres pruebas principales que validan los siguientes aspectos: 
\begin{description}
    \item Funcionamiento del formulario de login y registro.
    \item Búsqueda de videojuegos desde la barra de búsqueda.
    \item Creación de listas de videojuegos por parte del usuario.
\end{description}

Estas pruebas permiten comprobar que los componentes de la interfaz responden correctamente a interacciones del usuario y que la comunicación con el backend se realiza correctamente.

Cada prueba contiene:
\begin{enumerate}
    \item Verificación de la interfaz: se comprueba que los elementos de la interfaz existan y que se muestren correctamente.
    \item Simulación de la interacción del usuario: se simulan acciones como hacer clic en botones, rellenar formularios y navegar por la aplicación.
\end{enumerate}

En la figura \ref{FIG:CYPRESS} se muestra un ejemplo de prueba donde se comprueba que los elementos de la interfaz exiten y las respuestas del servidor son correctas 
y en la figura \ref{FIG:CYPRESS2} realiza una interacción del usuario.

\begin{figure}[Pruebas Cypress comprobación de elementos]{FIG:CYPRESS}{Salida de una prueba de Cypress que comprueba elementos de la interfaz y respuestas del servidor}
    \image{13cm}{}{cypress}
\end{figure}
\begin{figure}[Pruebas Cypress interacción del usuario]{FIG:CYPRESS2}{Salida de una prueba de Cypress que simula interacciones del usuario pinchando en un botón y rellenando un formulario}
    \image{13cm}{}{cypress2}
\end{figure}